\documentclass[11pt]{article}
\usepackage[margin=1in]{geometry}
\usepackage{mathpazo}
\usepackage{amsmath,amssymb,amsthm,mathtools,stmaryrd}
\usepackage[expansion=false]{microtype}
\usepackage{xcolor}
\usepackage{booktabs,tabularx,array,longtable}
\usepackage{enumitem}
\usepackage{listings}
\usepackage{tikz}
\usetikzlibrary{positioning,arrows.meta,fit,calc,backgrounds}
\usepackage{xurl}
\usepackage[hidelinks]{hyperref}
\usepackage{pdflscape}

\theoremstyle{definition}
\newtheorem{requirement}{Requirement}
\newtheorem{principle}{Principle}
\newtheorem{decision}{Open decision}
\newtheorem{obligation}{Obligation}
\theoremstyle{remark}
\newtheorem{remark}{Remark}

\newcommand{\fl}{f1r3lang}
\newcommand{\Camp}{CampF1R3}
\newcommand{\Kind}{K1ndl1ng}
\newcommand{\MUST}{\textsc{must}}
\newcommand{\MUSTNOT}{\textsc{must not}}
\newcommand{\SHOULD}{\textsc{should}}
\newcommand{\MAY}{\textsc{may}}
\usepackage[htt]{hyphenat}
\newcommand{\code}[1]{\texttt{#1}}
\setlength{\emergencystretch}{3em}

\definecolor{kw}{RGB}{0,0,140}
\definecolor{cm}{RGB}{90,110,90}
\lstdefinelanguage{Rho}{
  keywords={new,in,for,Nil,where,match,if,else,true,false,contract},
  sensitive=true,
  comment=[l]{//},
  morestring=[b]",
}
\lstset{
  language=Rho,
  basicstyle=\ttfamily\small,
  keywordstyle=\color{kw},
  commentstyle=\color{cm}\itshape,
  columns=fullflexible,
  keepspaces=true,
  frame=single,
  framerule=0.3pt,
  rulecolor=\color{black!30},
  xleftmargin=4pt,xrightmargin=4pt,
  aboveskip=6pt,belowskip=6pt,
  breaklines=true,
  literate={<-}{{$\leftarrow$}}1 {<=}{{$\Leftarrow$}}1
}

\newcommand{\Gz}{F1R3Gaze}
\newcommand{\enc}[1]{\llbracket #1 \rrbracket}
\newtheorem{workpackage}{Work package}
\newtheorem{finding}{Finding}
\lstdefinelanguage{RustLite}{
  keywords={pub,fn,struct,enum,trait,impl,for,type,let,mut,self,Self,use,mod,const,where,match,if,else,return,Some,None,Ok,Err,dyn,Box,Vec,Option,Result,u8,u16,u32,u64,i64,u128,bool,String,str},
  sensitive=true, comment=[l]{//}, morestring=[b]",
}
\title{\textbf{\Gz{}: Implementation Specification}\\[4pt]
\large A browser whose only execution mechanism is \fl{} on a native RSpace\\[6pt]
\normalsize Version 0.1 (draft for review), companion to the architecture proposal v0.2}
\author{Lucius Gregory Meredith\\
\small CEO, F1R3FLY.io, and Chief Science Officer, F1R3FLY Industries}
\date{28 September 2026}

\begin{document}
\maketitle

\begin{abstract}
\noindent
The architecture proposal in \code{publications/f1r3gaze} settles \emph{what} the
browser is: Blitz as the document engine, the \Camp{} executive as the only
execution mechanism, a F1R3FLY shard as the server, Servo quarantined for legacy
pages, and \code{f1r3x-wasm} as the reach tier. This document settles \emph{how}
to build it. It pins the code it builds on, records what a line-by-line review of
that code found --- including four places where the proposal assumes machinery
that does not yet exist --- and then specifies, to the level of types, wire forms
and algorithms, the upstream changes to \Camp{}, the \code{.knf} container, the
tab executive and its frame loop, the replay log, the DOM protocol, every
capability's protocol, the shard bridge and its light-client check, key-bound
sessions, the graded resolver, the process model, devtools, the reach tier, the
legacy tab, the projection gateway, and the conformance suite. It closes with work
packages mapped onto the proposal's five gates, and with the obligations and
decisions that remain open. The working name \Gz{} follows the name of the
directory in which the proposal lives.
\end{abstract}

\tableofcontents

%======================================================================
\section{Scope and conventions}\label{sec:scope}

\subsection{What this document is}
The proposal (v0.2) is normative for architecture; this specification is normative
for implementation. Where the two disagree, this document says so explicitly and
gives the reason, which in every case is something found in the code. The key words
\MUST{}, \MUSTNOT{}, \SHOULD{} and \MAY{} have their usual meaning.

\subsection{Pinned baselines}
Everything below was read at the commits in Table~\ref{tab:pins}. Work packages
\MUST{} re-pin at the start of each phase and re-run the findings of
Section~\ref{sec:findings} against the new pins.

\begin{table}[!htb]
\centering\small
\begin{tabularx}{\linewidth}{@{}lp{3.1cm}X@{}}
\toprule
Repository & Commit (date) & Role here\\
\midrule
\code{F1R3FLY-io/campf1r3} & \code{b46e16c} (18 Sep 2026) & the executive: 11 crates, $\approx$12{,}600 lines, 170 tests, no external dependencies\\
\code{F1R3FLY-io/f1r3node-rust} & \code{fce422a} (18 Sep 2026), \code{master} & the shard: HTTP API, \code{/ws/events}, \code{VersionedRegistry.rho}, \code{rho-pure-eval}, cost accounting\\
\code{F1R3FLY-io/publications} & HEAD, 28 Sep 2026 & \code{f1r3gaze} (proposal), \code{campf1r3} (store, kernel, bytecode specs), \code{fuzzyware} (graded where-clauses)\\
\code{F1R3FLY-io/embers} & HEAD, 28 Sep 2026 & prepare/sign/send deploy pattern; \code{firefly-client} crate\\
\code{F1R3FLY-io/F1R3Drive} & \code{0f3849a} & chunked on-chain file layout; write-back cache with bulk deploys\\
\code{F1R3FLY-io/F1R3Sky} & \code{57dcae5} & client-held keys through \code{@f1r3fly-io/embers-client-sdk}\\
\code{DioxusLabs/blitz} & \code{674d7d2} (28 Sep 2026), 0.3.0-beta.2 & document engine; Rust 1.91, edition 2024\\
\bottomrule
\end{tabularx}
\caption{Baselines. The node pins nightly-2026-02-09; \Camp{} targets Rust 1.75 and
edition 2021; Blitz requires 1.91 and edition 2024.}
\label{tab:pins}
\end{table}

\begin{requirement}[Toolchain]
The \Gz{} workspace \MUST{} build on stable Rust $\geq$1.91, edition 2024. \Camp{}
crates \MUST{} be consumed unmodified as git dependencies at a pinned revision
(edition-2021 crates compile inside an edition-2024 workspace), and every change to
them \MUST{} land upstream first (Section~\ref{sec:upstream}). No crate in the page
path \MAY{} require nightly.
\end{requirement}

%======================================================================
\section{Findings from the code review}\label{sec:findings}

Each finding states what the proposal assumes, what the code does, and what this
specification does about it.

\begin{finding}[No ground data in the kernel]\label{f:ground}
The proposal's examples send strings (\code{doc!("query1", "\#lamp", *found)}). The
\Kind{} specification lists ``ground literals of every sort, collections, methods''
as explicitly out of scope, and \code{k1ndl1ng-norm}'s \code{Node} has no literal
constructor. \emph{Response:} a ground-data sort, K1G, specified in
\S\ref{sec:u-ground} and scheduled into P0.
\end{finding}

\begin{finding}[Guards are parsed and discarded]\label{f:guards}
At K2 the parser recognises \code{where} and then skips tokens to the closing
parenthesis; nothing is recorded and nothing evaluates. The node's
\code{rho-pure-eval} cannot be reused in the client: it depends on \code{models},
which depends on \code{rspace++} (\code{heed}, \code{tokio}) and \code{tonic}.
\emph{Response:} a portable evaluator, \code{k1ndl1ng-eval}
(\S\ref{sec:u-guards}), differentially tested against \code{rho-pure-eval}.
\end{finding}

\begin{finding}[No host channels, and no way to meter]\label{f:apertures}
\code{b3ll0ws::Machine} owns an \code{RSpace} fixed to \code{NoObserver}, so the
page meter cannot be installed; and a send on a grounded name can only be seen by
the host by polling the store. \emph{Response:} a \code{Machine} generic over its
observer, and \emph{apertures} (\S\ref{sec:u-apertures}).
\end{finding}

\begin{finding}[No oracle seam]\label{f:oracle}
The proposal speaks of an oracle $\omega$ to which a Graded profile attaches. In the
code, choice is first-match inside \code{Space::try\_candidate}, and \code{Profile}
is an enum whose \code{Free} variant is currently treated as \code{Deterministic}.
\emph{Response:} a \code{Resolver} trait over enumerated candidates
(\S\ref{sec:u-resolver}).
\end{finding}

\begin{finding}[Mixed persistence in joins is well defined here]
\code{b3ll0ws::receipt\_parts} makes a continuation persistent if \emph{any} bind is
persistent, while data persistence stays per datum. The proposal's lamp,
\code{for (\_ <= clicks \& \_ <- off)}, therefore means ``a persistent handler that
consumes one click and the \code{off} token'', which is what was intended. The
example stands; Appendix~\ref{app:lamp} gives it in K1G form.
\end{finding}

\begin{finding}[Blitz dispatch is synchronous]\label{f:sync}
\code{blitz\_dom::EventHandler::handle\_event} receives an \code{EventState} and must
decide \code{prevent\_default} and \code{stop\_propagation} before it returns. The
proposal's ``apply the default after a deadline'' cannot be implemented on it.
\emph{Response:} static listener options, plus an opt-in synchronous drain bounded
by a count of communications (\S\ref{sec:events}).
\end{finding}

\begin{finding}[Blitz is one process, and already hosts a script engine]\label{f:blitz}
\code{blitz-shell} drives a \code{Box<dyn Document>} in the window's process;
there is no renderer/browser split. \code{anyrender\_serialize} exists and makes a
display-list transport feasible. Blitz now also ships \code{blitz-vibey-script}, a
Boa-backed \code{ScriptDocument}; its shape (wrap \code{BaseDocument}, implement
\code{Document::handle\_ui\_event} and \code{Document::poll}) is the template for
\code{RhoDocument}. \emph{Response:} one OS process with a worker thread per tab
through P3; process separation in P4 (\S\ref{sec:process}). \code{blitz-vibey-script}
\MUSTNOT{} be linked into any binary that hosts \fl{} pages.
\end{finding}

\begin{finding}[Node API, registry and costs]\label{f:node}
The node serves more than the proposal lists: \code{/api/explore-deploy},
\code{/api/explore-deploy-by-block-hash}, \code{/api/data-at-name-by-block-hash},
\code{/api/registry/\{uri\}} with a \code{block\_hash} query parameter,
\code{/api/estimate-cost}, \code{/api/prepare-deploy},
\code{/api/deploy-finalization-status/\{sig\}}, and \code{/ws/events} carrying
\code{block-added}, \code{block-created}, \code{block-finalised} and deploy
lifecycle events. There is no inclusion-proof endpoint. In
\code{VersionedRegistry.rho} the header records that Step~6 (the public
\code{rho:registry:1.0.0} entry point) and Step~7 (notify firing) were pending; the
notify list is captured but ``ignored for now''. Storage cost is computed from
\code{prost}-encoded sizes of \code{Par}s. \emph{Response:} freshness by
re-resolution on \code{block-finalised} until Step~7 lands; the client meter is a
local budget, and deploy costs are quoted from \code{/api/estimate-cost}
(\S\ref{sec:meter}).
\end{finding}

\begin{finding}[Session names in \code{f1r3x-serve} are bearer tokens]\label{f:serve}
A client attaches at a path of the form
\code{ws://host/s/<session-name>}, choosing the name itself,
and ``an unknown name creates the session''. Whoever learns the name owns the
session. \emph{Response:} authenticated attach (\S\ref{sec:u-serve}).
\end{finding}

\begin{finding}[Existing F1R3FLY clients already do half the shard bridge]\label{f:embers}
Embers exposes each mutation as \code{/prepare} (the server renders a \fl{} template
and returns an unsigned contract under a JWT) and \code{/send} (the client returns it
signed); F1R3Sky holds the private key in the client through the Embers SDK.
Embers' \code{firefly-client} crate already signs \code{DeployData} with
\code{secp256k1}, submits over gRPC and HTTP, and consumes \code{/ws/events}.
F1R3Drive stores a file as a metadata map with a first chunk inline and further
chunks on sub-channels, deployed in bulk from a write-back cache.
\emph{Response:} the shard bridge depends on \code{firefly-client}
(\S\ref{sec:shard}); F1R3Drive's layout is the on-chain fallback for blobs
(\S\ref{sec:blobs}).
\end{finding}

\begin{finding}[One missing adapter]
The proposal cites a \code{campf1r3-ispace} adapter for differential replay against
\code{rspace++}. It is not in the pinned repository. Obligation~\ref{ob2:replay}
carries it.
\end{finding}

%======================================================================
\section{Workspace}\label{sec:workspace}

A new repository, \code{F1R3FLY-io/f1r3gaze}, Apache-2.0, holds everything that is
not an upstream change. Crates, in dependency order:

\begin{center}\small
\begin{tabularx}{\linewidth}{@{}p{3.3cm}p{1.5cm}X@{}}
\toprule
Crate & Target & Contents\\
\midrule
\code{gaze-ground} & all & \code{Ground} values, conversion to and from K1G normal forms, the reply-tuple conventions of \S\ref{sec:proto-conv}\\
\code{gaze-knf} & all & the \code{.knf} container: reader, writer, manifest, integrity strings\\
\code{gaze-exec} & all & \code{TabExec}: executive, apertures, meter, resolver, recorder, frame loop\\
\code{gaze-graded} & all & semirings, resolvers, deterministic PRNG, fairness, ceiling check\\
\code{gaze-dom-core} & all & the DOM protocol engine over a \code{DomBackend} trait; name table; attenuation; event routing\\
\code{gaze-dom-blitz} & native & \code{DomBackend} over \code{blitz-dom}; \code{RhoDocument}; \code{RhoEventHandler}\\
\code{gaze-broker} & native & manifests to grants, policy, prompts, routing table, revocation\\
\code{gaze-net} & native & \code{blitz\_traits::net::NetProvider} with hash verification; the \code{net} capability\\
\code{gaze-store} & native & per-origin key-value store behind the \code{store} capability\\
\code{gaze-shard} & native & shard bridge: keystore, deploys, registry, reads, proofs, sessions, watches\\
\code{gaze-blob} & native & \code{BlobSource} implementations and the content cache\\
\code{gaze-shell} & native & the application: windows, tabs, chrome, prompts, settings\\
\code{gaze-devtools} & all & debug capability; the devtools page is itself \fl{}\\
\code{gaze-reach} & wasm32 & \code{cdylib} bundling the executive and \code{gaze-dom-core}, plus \code{gaze-reach.js}\\
\code{gaze-legacy} & native & P4: Servo in its own process\\
\code{gaze-gateway} & native & P4: headless projection server\\
\code{gaze-conformance} & all & the protocol and determinism suites of \S\ref{sec:testing}\\
\code{f1r3c} & native & toolchain CLI: surface \fl{} or kernel text to \code{.knf}; site publishing\\
\bottomrule
\end{tabularx}
\end{center}

\begin{requirement}[Portable core]
\code{gaze-ground}, \code{gaze-knf}, \code{gaze-exec}, \code{gaze-graded},
\code{gaze-dom-core} and \code{gaze-devtools} \MUST{} build for
\code{wasm32-unknown-unknown}, \MUST{} carry \code{\#![forbid(unsafe\_code)]}, and
\MUSTNOT{} depend on \code{std::time}, threads, sockets, floating point on any path
that affects output, or any crate outside the Rust standard library and \Camp{}.
\end{requirement}

\noindent
This is the rule that makes the native browser and the reach tier one
implementation: the protocol engine and the executive are the same code in both,
and only the \code{DomBackend} differs.
%======================================================================
\section{Upstream changes to \Camp{}}\label{sec:upstream}

Every change here is application-agnostic, so it respects the constraint recorded
when \Camp{} was split from Rhobots-in-Space: nothing in that repository may know
that a browser exists. Each change is behind a cargo feature, off by default, so the
Rhobots path and the 170 existing tests are unaffected.

\subsection{U1: ground data (level K1G, feature \code{ground})}\label{sec:u-ground}

\paragraph{Sort.} A new process constructor carries ground values:
\begin{lstlisting}[language=RustLite]
pub enum Node { /* existing variants */, Ground(Lit) }
pub enum Lit {
    Bool(bool), Int(i64), Str(Box<str>), Bytes(Box<[u8]>),
    List(Vec<Norm>), Tuple(Vec<Norm>), Map(Vec<(Norm, Norm)>),
}
\end{lstlisting}
Collection elements are arbitrary normal forms, so names (as \code{*x}) may sit
inside lists, tuples and maps. There are no floats: every consumer in this
specification needs bit-identical results across targets.

\paragraph{Encoding.} Tags continue the table of the \Kind{} specification in the
unused range between \code{Wild} (\code{0x08}) and \code{Quote} (\code{0x10}):

\begin{center}\small
\begin{tabular}{@{}lll@{}}
\toprule
Tag & Constructor & Payload\\
\midrule
\code{0x09} & \code{Bool} & one byte, \code{0} or \code{1}\\
\code{0x0A} & \code{Int} & zigzag LEB128 \code{i64}\\
\code{0x0B} & \code{Str} & LEB128 length, UTF-8 bytes (no normalisation)\\
\code{0x0C} & \code{Bytes} & LEB128 length, bytes\\
\code{0x0D} & \code{List} & $n$, $\enc{e_1}\cdots\enc{e_n}$, source order\\
\code{0x0E} & \code{Tuple} & $n$, $\enc{e_1}\cdots\enc{e_n}$, source order\\
\code{0x0F} & \code{Map} & $n$, then pairs sorted by $\enc{k}$ bytewise; duplicate keys are a normalisation error\\
\bottomrule
\end{tabular}
\end{center}

\noindent
Requirement \emph{order key is the encoding itself} of the \Kind{} specification is
preserved. Whether bytewise order on these tags agrees with the node's
\code{models} \code{Score} order for the corresponding \code{Expr}s is added to that
specification's sort-agreement obligation; until it is discharged, a normal form
containing ground data \MUSTNOT{} be claimed equal to a node term by hash.

\paragraph{Surface syntax.} Rholang's literal syntax, as \code{rholang-tree-sitter}
fixes it: \code{true}, \code{false}, decimal integers with optional sign, double-quoted
strings with rholang's escapes, \code{[\ldots]}, \code{(a, b, \ldots)} with at least
two elements, \code{\{k: v, \ldots\}}. Bytes have no surface form at K1G; they occur
only in \code{.knf} encodings and in data injected by a host, and the printer shows
them as \code{0x\ldots}. Methods remain out of scope.

\paragraph{Matching.} In \code{k1ndl1ng-campf1r3::spatial}, a ground pattern matches
a ground datum iff their encodings are equal; collection patterns match elementwise,
with \code{FreeVar}, \code{Wild} and \code{Eval(Name::Free)} binding inside them as
they already do inside processes. Map patterns match exactly (no row
variables) in v1.

\subsection{U2: guards (level K2, feature \code{guards})}\label{sec:u-guards}

A new crate, \code{k1ndl1ng-eval}, pure and portable, evaluates the guard subset:
ground literals, bound variables, \code{== != < <= > >=} on \code{Int} and on
\code{Str} (bytewise), \code{+ - * / \%} on \code{Int} with checked arithmetic
(overflow and division by zero make the guard evaluate to an error, which is
\emph{not enabled}), \code{++} on \code{Str}, \code{Bytes} and \code{List},
\code{and}, \code{or}, \code{not}, and two graded scalars used only in graded pages:
\code{w(n, d)} and \code{cost(n)} (\S\ref{sec:graded}). The parser records the guard
as a tree instead of skipping it; the normal form stores it on the \code{Receive}
node under a new tag \code{0x21} (\code{Guard}), after the binds and before the body.

\begin{requirement}[Guard evaluation site]
Guards \MUST{} be evaluated in \code{Matcher::check\_commit(k, bound)}, which is
where the store already asks, once per candidate. For crisp pages the result is a
Boolean; for graded pages it is routed to the resolver (U5).
\end{requirement}

\subsection{U3: an observer-generic machine}\label{sec:u-observer}
\begin{lstlisting}[language=RustLite]
pub struct Machine<B: Backend = Direct, M: NameMinter = Sequential,
                   O: Observer<Chan, BindPat, Data, Cont> = NoObserver> { .. }
impl<B, M, O> Machine<B, M, O> {
    pub fn with_observer(cfg: Config, backend: B, minter: M, observer: O) -> Self;
}
\end{lstlisting}
\code{k1ndl1ng-campf1r3::rspace} gains \code{rspace\_with(cfg, observer)}.
\code{f1r3x::Executive} keeps its signature.

\subsection{U4: apertures}\label{sec:u-apertures}

An aperture is a channel owned by the host. The page can send on it; the store
never sees those sends.

\begin{lstlisting}[language=RustLite]
pub struct ApertureId(pub u32);
pub struct Outbound { pub aperture: ApertureId, pub args: Vec<Norm>, pub seq: u64 }

impl<B, M, O> Machine<B, M, O> {
    /// Sends on `chan` go to the outbox instead of the store.
    pub fn bind_aperture(&mut self, chan: Name, id: ApertureId);
    pub fn unbind_aperture(&mut self, id: ApertureId);
    pub fn take_outbound(&mut self) -> Vec<Outbound>;
    /// Host-originated send: enqueued as a `Send` job at the tail of the queue.
    pub fn inject(&mut self, chan: Name, args: Vec<Norm>);
}
\end{lstlisting}

\begin{itemize}[nosep]
\item A \code{Send} job whose channel is bound to an aperture is appended to the
  outbox with a monotone \code{seq}, emits \code{Step::Outbound}, and calls a new
  \code{Observer::observe\_aperture} hook, which may abort like any other.
\item A persistent send on an aperture is delivered once and emits a diagnostic.
\item A receipt on an aperture channel is stored and is inert: the host never
  produces on apertures, it injects on names the page gave it.
\item \code{inject} is the only host write path. It does not bypass the store,
  the observer or the resolver.
\end{itemize}

\subsection{U5: the resolver seam and the Graded profile}\label{sec:u-resolver}

\code{Profile} gains \code{Graded}. In that profile \code{try\_candidate} does not
return the first admissible candidate; it enumerates up to
\code{Config::max\_candidates} (default 256) admissible candidates in the existing
deterministic order, asks the matcher for each one's value, and asks the resolver to
choose.

\begin{lstlisting}[language=RustLite]
pub trait Matcher<P, A, K> {
    fn get(&self, p: &P, a: &A) -> Option<A>;
    fn check_commit(&self, k: &K, bound: &[A]) -> bool { true }
    /// Graded profile only: the clause's value, opaque to the store.
    fn weigh(&self, k: &K, bound: &[A]) -> Weight { Weight::ONE }
}
pub struct Weight(pub u128);            // semiring-specific packing
pub trait Resolver: Send + Sync + 'static {
    /// `ws` in candidate order; None means no candidate is enabled.
    fn select(&self, site: Hash32, ws: &[Weight]) -> Option<usize>;
}
\end{lstlisting}

\noindent
The log gains \code{Event::Select \{ site, n, chosen \}}. In \code{Replay} the
selection is read back from the log rather than recomputed, so sampled runs replay
exactly. Truncation at \code{max\_candidates} emits a diagnostic and is recorded.

\begin{remark}[Local versus contention-set resolution]\label{rem:local}
The store resolves at a cut: the candidates are those that use the arriving record.
The quantum paper normalises over the \emph{contention component}, the connected
component of the conflict relation over the whole term. For idempotent semirings
(Boolean, Viterbi, tropical) the cut-local argmax is what a page asks for. For
sampled $\mathbb{R}_{\geq 0}$ pages it is an approximation of the paper's
semantics, and the specification says so rather than hiding it. An exact
\emph{round mode}, which collects the enabled sites of an admitted round, builds the
conflict graph on shared occurrences and resolves per component, is
Obligation~\ref{ob2:round}.
\end{remark}

\subsection{U6: a keyed minter}
\code{Sequential} is a pure function of binder path and counter, which is right for
Rhobots and wrong where names might be exported. \code{Keyed \{ seed: [u8; 32],
counter: u64 \}} mints \code{blake2b\_256("k1ndl1ng/keyed/" || seed || counter ||
path)}. The browser seeds it from the OS at tab creation and records the seed in
the replay log.

\subsection{U7: authenticated attach in \code{f1r3x-serve}}\label{sec:u-serve}
\code{ready} gains a 32-byte nonce. A client attaches with
\code{auth <pubkey-hex> <sig-hex>} over the nonce and the session path, and the
session is keyed by the public key; the path segment becomes a label. An
unauthenticated attach is refused unless the server is started with
\code{--open}, which the browser never uses. The Rhobots client is updated in the
same change.

%======================================================================
\section{The \code{.knf} container and how pages load}\label{sec:knf}

\subsection{Format}
\begin{center}\small
\begin{tabularx}{\linewidth}{@{}lX@{}}
\toprule
Field & Contents\\
\midrule
magic & \code{F1R3KNF\textbackslash0} (8 bytes)\\
version & \code{u16}, 1\\
level & \code{u8}: 0 K0, 1 K1, 2 K1G, 3 K2 (implies K1G)\\
sections & repeated \code{(tag u8, len LEB128, bytes)}, tags ascending, each at most once\\
\quad\code{0x01} manifest & a K1G \code{Map} (below), encoded\\
\quad\code{0x02} body & the \emph{ungrounded} normal-form encoding\\
\quad\code{0x03} source map & optional; spans for diagnostics and devtools\\
\bottomrule
\end{tabularx}
\end{center}

\noindent
The manifest is:
\begin{lstlisting}
{ "imports":  [ ("doc", "rho:gaze:doc"), ("net", "rho:gaze:net"), ... ],
  "semiring": "boolean" | "viterbi" | "tropical" | "prob",
  "fair":     false,
  "ceiling":  "crisp" | "idempotent" | "sampled",
  "budget":   { "frame": 20000, "sync": 512 } }
\end{lstlisting}
\code{imports} lists every free identifier of the body in level order with the
capability it requests. An identifier the manifest does not list is a load error.

\subsection{Identity}
\begin{requirement}[Two hashes]
The \emph{program hash} is \code{k1ndl1ng-norm}'s BLAKE2b-256 of the body, exactly
as the executive computes it; it keys the parse and plan caches. The \emph{grant
hash} is BLAKE2b-256 of \code{program hash || enc(manifest)}; it keys stored grants,
because the manifest, not the body, states the authority requested.
The HTML \code{integrity} attribute \MUST{} carry the program hash as
\code{blake2b-256:<64 hex>}.
\end{requirement}

\subsection{In HTML}
\begin{lstlisting}
<script type="application/f1r3lang" src="app.knf"
        integrity="blake2b-256:9f2c...e1"></script>
\end{lstlisting}
Inline kernel text is accepted with \code{level}, \code{imports} and
\code{semiring} attributes standing in for the manifest; the loader parses it at the
declared level and builds the manifest from the attributes. A document may carry
several \fl{} scripts; each is a separate deploy into the same tab executive, in
document order, sharing grants by identifier. Scripts of any other type are never
executed; if a document has JavaScript scripts and no \fl{} script, the shell offers
to reopen it in a legacy tab (\S\ref{sec:legacy}).

\subsection{Load sequence}
\begin{enumerate}[nosep]
\item Blitz parses the document; \code{RhoDocument} collects \fl{} scripts.
\item Each \code{src} is fetched through \code{gaze-net}; its program hash is
  recomputed and compared with \code{integrity}. A mismatch is fatal to the page.
\item The broker turns the manifest into grants (\S\ref{sec:broker}), prompting where
  policy requires. The tab shows the grants before anything runs.
\item For each import the broker returns a 32-byte key; the executive binds an
  aperture on \code{Name::Unforgeable(key)} and grounds the body with
  \code{ground\_free(keys)}. Denied imports are grounded to fresh keys that are
  bound to nothing: dead channels.
\item The grounded term is spawned; the first frame runs.
\end{enumerate}

%======================================================================
\section{The tab executive}\label{sec:tabexec}

\subsection{Structure}
\begin{lstlisting}[language=RustLite]
pub struct TabExec {
    machine:   Machine<Direct, Keyed, PageMeter>,
    apertures: BTreeMap<ApertureId, Service>,  // doc, net, store, ...
    inbox:     Vec<Injection>,                 // filled by services between frames
    recorder:  Recorder,                        // the replay log, section 6.3
    dom:       gaze_dom_core::Engine,
    frame:     u64,
}
\end{lstlisting}

\subsection{The frame}
\code{TabExec::frame(now\_ms)} is the whole driver. It is a pure function of the
executive state, the inbox and \code{now\_ms}.

\begin{enumerate}[nosep]
\item \textbf{Order the inbox} by (source class, arrival sequence) with classes in
  the fixed order: DOM replies, timers, input events, network, store, shard. Record
  \code{Frame\{n, now\_ms\}} and then each \code{Inject}.
\item \textbf{Inject} each entry with \code{Machine::inject}.
\item \textbf{Run} \code{machine.run(budget)} where \code{budget} is the manifest's
  frame budget in steps, reduced by the meter's remaining allowance.
\item \textbf{Drain} \code{take\_outbound()} in \code{seq} order. DOM messages go to
  \code{gaze-dom-core}, which answers reads from the snapshot taken at the start of
  the frame and appends writes to the frame's commit batch; everything else goes to
  its service, which will answer by adding to the inbox of a later frame.
\item \textbf{Commit}: the batch is handed to the \code{DomBackend} and applied with
  one \code{DocumentMutator} and one \code{flush()}. Replies the DOM engine produced
  (names of created nodes, read results) enter the inbox for frame $n+1$.
\end{enumerate}

\begin{requirement}[Frame atomicity, restated]
Writes produced in frame $n$ \MUST{} become visible to layout together, after step
5. Reads in frame $n$ \MUST{} observe the document as committed at the end of frame
$n-1$, and their replies \MUST{} arrive in frame $n+1$. A page therefore never
observes a half-applied frame, and a read is never answered from state that another
read in the same frame could not see.
\end{requirement}

\begin{remark}[Why not \code{Session::poll}]
The proposal drives the executive through \code{f1r3x-host::Session::poll}. A
\code{Session} exists to serialise rounds and scenes for a client across a boundary;
in the page process the renderer is in-process and needs the outbox, not scenes.
\code{TabExec} therefore drives the machine directly and keeps the host's admission
rule for the debugger's lockstep mode. \code{Session} is used where it belongs: when
devtools or a Rhobots headset attaches to a tab (\S\ref{sec:devtools}).
\end{remark}

\subsection{The replay log}\label{sec:log}
A \code{.gzlog} file is a header and a record stream.
\begin{description}[nosep,leftmargin=1.5em,style=nextline]
\item[Header] program hashes of all scripts, grant hash, grants by capability kind
  (never keys or secrets), minter seed, profile, semiring, executive and bridge
  build identifiers.
\item[Records] \code{Frame\{n, now\_ms\}}; \code{Inject\{class, aperture\_or\_name,
  args\}} with arguments as K1G encodings; \code{Sync\{listener, event\}} for
  synchronous drains; \code{Select\{site, n, chosen\}} copied from the store log;
  \code{Commit\{n, batch\_hash\}}; \code{Budget\{n, exhausted\}}.
\end{description}
Replay loads the header, rebuilds the tab with the same seed and with services
replaced by the log, and re-injects each record at its frame. The observable result
--- the committed document at every frame, compared by \code{batch\_hash} --- \MUST{}
be identical on \code{x86\_64}, \code{aarch64} and \code{wasm32}.

%======================================================================
\section{The DOM protocol}\label{sec:proto}

This section discharges the proposal's Obligation on the DOM protocol. It is
implemented once, in \code{gaze-dom-core}, over:
\begin{lstlisting}[language=RustLite]
pub trait DomBackend {
    type Node: Copy + Ord;
    fn query_in(&self, root: Self::Node, sel: &str, all: bool) -> Result<Vec<Self::Node>, SelErr>;
    fn read(&self, n: Self::Node, what: Read) -> ReadResult;
    fn apply(&mut self, batch: &[Write<Self::Node>]) -> Vec<Created<Self::Node>>;
    fn is_descendant(&self, n: Self::Node, of: Self::Node) -> bool;
}
\end{lstlisting}
\code{gaze-dom-blitz} implements it with \code{BaseDocument::query\_selector\_all\_in},
\code{DocumentMutator} and \code{NodeId}; \code{gaze-reach} implements it over the
host browser's DOM.

\subsection{Conventions}\label{sec:proto-conv}
\begin{itemize}[nosep]
\item Every request is a send on an element name: \code{E!(verb, args\ldots, ret)}
  where \code{verb} is a \code{Str} and \code{ret} is a name, or \code{E!(verb,
  args\ldots)} for writes that need no acknowledgement.
\item Every reply is one tuple on \code{ret}: \code{("ok", v)} or
  \code{("err", code, detail)}. Receivers match with patterns such as
  \code{for (@("ok", *el) <- r)}.
\item Element names are unforgeable names minted by the engine, one per live node,
  stable for the node's life: asking twice for the same node yields the same name,
  so name equality is node equality.
\item Error codes: \code{"attenuated"} (outside the holder's subtree),
  \code{"detached"} (node removed), \code{"selector"}, \code{"type"} (wrong argument
  sort), \code{"verb"}, \code{"quota"}, \code{"revoked"}.
\end{itemize}

\begin{principle}[Attenuation is the tree, enforced]
Every verb that returns or addresses a node \MUST{} check
\code{is\_descendant(target, holder\_root)} where \code{holder\_root} is the node of
the name the request was sent on. \code{parent} and \code{closest} never return a node
above the holder.
\end{principle}

\subsection{Verbs}
\begin{center}\small
\begin{tabularx}{\linewidth}{@{}p{4.4cm}X@{}}
\toprule
Request & Meaning and reply\\
\midrule
\code{query1(sel, ret)} & first match in the subtree; \code{("ok", *E)} or \code{("err","none",sel)}\\
\code{query(sel, ret)} & all matches; \code{("ok", [*E1, ...])}\\
\code{children(ret)}, \code{parent(ret)}, \code{closest(sel, ret)} & navigation, attenuated\\
\code{attr(k, ret)}, \code{setAttr(k, v)}, \code{removeAttr(k)} & attributes; \code{attr} replies \code{("ok", v)} or \code{("ok", Nil)}\\
\code{text(ret)}, \code{setText(s)} & text content\\
\code{value(ret)}, \code{setValue(s)} & form-control values\\
\code{classAdd(c)}, \code{classRemove(c)}, \code{classToggle(c)} & class list\\
\code{style(k, v)}, \code{unstyle(k)} & inline style properties\\
\code{append(frag, ret)}, \code{prepend}, \code{before}, \code{after}, \code{replaceWith} & insert a fragment; replies \code{("ok", \{ref: *E\})} for every \code{("ref", k)} marker\\
\code{setHTML(html)} & parse with \code{html5ever} and replace children; scripts in the fragment are inert\\
\code{remove()} & detach; the name dies\\
\code{rect(ret)} & layout box as \code{(x, y, w, h)} in integer CSS pixels, from the last committed layout\\
\code{focus()}, \code{blur()}, \code{scrollIntoView()} & focus and scroll\\
\code{listen(type, ch, opts, ret)} & subscribe (\S\ref{sec:events}); replies \code{("ok", *sub)}; \code{sub!("cancel")}\\
\bottomrule
\end{tabularx}
\end{center}

\noindent
A fragment is either a \code{Str} of HTML or a structured term, which avoids
string parsing and is the preferred form:
\begin{lstlisting}
("el", "li", {"class": "todo"}, [ ("text", "buy milk"), ("ref", "check") ])
\end{lstlisting}
The document root name \code{doc} additionally answers \code{title}, \code{setTitle},
\code{head}, \code{body} and \code{location}. Navigation is the \code{nav}
capability, not a document verb.

\subsection{Events}\label{sec:events}

\code{listen(type, ch, opts, ret)} with \code{opts} a map of Booleans:
\code{capture}, \code{once}, \code{prevent}, \code{stop}, \code{decide}.
Each occurrence is injected as one datum on \code{ch}:
\begin{lstlisting}
( type, { "target": *E, "x": 120, "y": 44, "button": 0, "key": "Enter",
          "mods": ["shift"], "value": "...", "reply": *R } )
\end{lstlisting}
Fields absent for a type are absent from the map. \code{target} is always inside
the listener's subtree, since the listener was installed on a name for that subtree.

\code{RhoEventHandler} implements \code{blitz\_dom::EventHandler}. For each
\code{DomEvent} it walks Blitz's chain, capture phase then bubble phase, and for
each node with listeners of that type:
\begin{enumerate}[nosep]
\item applies static options: \code{prevent} calls \code{EventState::prevent\_default},
  \code{stop} calls \code{stop\_propagation} after this node;
\item if \code{decide} is false (the default), queues the datum for the next frame;
\item if \code{decide} is true, injects the datum now and runs the tab executive
  synchronously for at most the manifest's \code{sync} budget in steps, then reads
  whatever \code{"prevent"} and \code{"stop"} messages arrived on \code{R}. The drain
  is recorded as \code{Sync}. If the budget expires, the default action proceeds.
\end{enumerate}
The budget is in steps, not milliseconds, so the decision is deterministic and
replayable. Text editing, caret movement and IME composition are Blitz's; the page
sees \code{input} events with \code{value}.
%======================================================================
\section{Capabilities and the broker}\label{sec:broker}

\subsection{Grants}
The broker maps each manifest import to a grant. A grant is a routing-table entry
\code{(tab, ApertureId) $\mapsto$ (service, attenuation, quota)}; the page holds only
the aperture's unforgeable name.

\begin{principle}[Revocation is a table update]
The proposal revokes by killing a forwarder process. With apertures the forwarder is
the routing entry: removing it makes every later send on the name answer
\code{("err","revoked")}, and the page cannot tell a revoked name from a dead one.
\end{principle}

\begin{center}\small
\begin{tabularx}{\linewidth}{@{}p{2.7cm}p{3.1cm}X@{}}
\toprule
Capability & Default policy & Attenuation\\
\midrule
\code{rho:gaze:doc} & granted & the document root; components get subtrees by name passing\\
\code{rho:gaze:log} & granted & console of this tab\\
\code{rho:gaze:clock} & granted & frame ticks, timers, frame-quantised time\\
\code{rho:gaze:rand} & granted & OS entropy, recorded in the log\\
\code{rho:gaze:net} & granted, same site & origin allow-list; widened only by prompt\\
\code{rho:gaze:store} & granted & this origin's space; quota\\
\code{rho:gaze:nav} & granted & same-site navigation; cross-site opens a new tab\\
\code{rho:gaze:shard} & prompt & per verb class: read, explore, deploy, session\\
\bottomrule
\end{tabularx}
\end{center}

\noindent
Grants are stored per (site identity, grant hash); a new manifest is a new request.
The shell lists a tab's grants and open sessions, and any entry can be revoked.

\subsection{Protocols}
\begin{description}[leftmargin=1.5em,style=nextline,itemsep=2pt]
\item[\code{clock}] \code{("frames", ch)} injects \code{("tick", n, now\_ms)} each
  frame until \code{("stop", ch)}; \code{("after", ms, ch)} fires once; \code{("now",
  ret)} replies with the current frame's time.
\item[\code{log}] \code{(level, v)} with level \code{"debug"}, \code{"info"},
  \code{"warn"} or \code{"error"}; \code{v} is printed with \code{show\_pretty}.
\item[\code{rand}] \code{("u64", ret)}, \code{("bytes", n, ret)}.
\item[\code{net}] \code{("fetch", \{"url", "method", "headers", "body",
  "integrity"\}, ret)} replies \code{("ok", \{"status", "headers", "body",
  "digest"\})}. With \code{integrity} set, the body is verified before delivery and a
  mismatch is \code{("err","integrity",url)}. Streaming and WebSockets are deferred;
  the shard bridge covers live data.
\item[\code{store}] \code{("get", k, ret)}, \code{("put", k, v, ack)}, \code{("del", k,
  ack)}, \code{("list", prefix, ret)}. Keys are \code{Str}; values are any closed
  normal form, stored as encodings. Unforgeable names inside stored values are
  rejected, because a name must not outlive the tab that minted it.
\item[\code{nav}] \code{("go", url)}, \code{("replace", url)}, \code{("back")},
  \code{("state", v)}.
\item[\code{shard}] \S\ref{sec:shard}.
\end{description}

\noindent
\code{gaze-store} keeps one append-only file per origin with periodic compaction.
The proposal's ``persistent space as a \Camp{} checkpoint'' is replaced by this
key-value protocol: it needs no upstream store serialisation and is simpler to
quota.

\subsection{The page meter}\label{sec:meter}
\code{PageMeter} implements \code{Observer}. It debits a per-tab allowance for each
produce, consume, COMM and aperture send, plus a per-byte charge on datum
encodings, from a table in \code{gaze-exec}. The allowance is refilled per frame,
larger for the foreground tab. Exhaustion aborts the operation atomically; the tab
stops making progress, the shell shows why, and the user can grant more. The meter
is a local budget and \MUSTNOT{} be presented as a prediction of shard cost:
shard cost is computed from \code{prost} sizes of \code{Par}s
(Finding~\ref{f:node}), so deploy costs are always quoted from
\code{/api/estimate-cost}.

%======================================================================
\section{The shard bridge}\label{sec:shard}

\subsection{Components}
\code{gaze-shard} runs in the browser process on \code{tokio} and depends on
Embers' \code{firefly-client} for \code{DeployData} construction, \code{secp256k1}
signing, submission and the \code{/ws/events} stream. It adds the keystore, the
registry resolver, the assurance ladder, the light client and sessions. Keys live in
the operating system's credential store; each (user, site) pair gets its own key,
generated on first grant, so sites cannot correlate a user by key.

\subsection{Addresses}
\begin{lstlisting}
f1r3://<publisher-pubkey-hex>/<project>@<semver-range>/<path>
f1r3h://blake2b-256/<64 hex>
\end{lstlisting}
The first resolves through \code{rho:serve:1:<pubkey>:<project>:<range>}; the second
names immutable content directly and may be fetched from anywhere.

\subsection{Site manifest}
A published version's registry entry holds, as its code, a process that sends a K1G
map:
\begin{lstlisting}
{ "gaze": 1, "entry": "index.html",
  "files": { "index.html": 0x9f2c..., "app.knf": 0x41aa..., "style.css": 0x07be... },
  "mirrors": [ "https://cdn.example.org/blob/" ] }
\end{lstlisting}
Every file is named by its BLAKE2b-256 hash, so the node never has to agree with
the client on how to hash a \code{Par}; this settles the proposal's obligation on
canonical code hashes for sites, because code travels as \code{.knf} bytes
addressed by hash rather than as registry terms.

\subsection{Resolution}
\begin{enumerate}[nosep]
\item Take the last finalized block from \code{/api/last-finalized-block} on each
  configured observer.
\item Query \code{/api/registry/\{uri\}?block\_hash=B} for the versioned URN.
\item Verify the answer at the configured rung (\S\ref{sec:ladder}).
\item Reject if \code{B} is older than the highest finalized block at which this
  publisher's binding was previously seen (freshness).
\item Fetch files by hash from any source (\S\ref{sec:blobs}); verify each hash.
\item Subscribe: on each \code{block-finalised} event, re-resolve open sites and
  notify tabs of changes. When \code{VersionedRegistry} Step~7 lands, subscribe to the
  notify channel instead.
\end{enumerate}

\subsection{The assurance ladder in code}\label{sec:ladder}
Every shard reply to a page carries its rung as the second tuple field, as in
\code{("ok", rung, v)}, where \code{rung} is one of \code{"node"}, \code{"quorum"},
\code{"proof"} or \code{"replayed"}.

\begin{description}[leftmargin=1.5em,style=nextline,itemsep=2pt]
\item[node] one observer's answer.
\item[quorum] identical canonical encodings from $k$ of $n$ configured observers
  run by distinct operators (default $k=2$, $n=3$).
\item[proof] requires node work package N1: \code{GET /api/proof/data-at-name} and
  \code{GET /api/proof/registry/\{uri\}}, each returning the block header,
  justification signatures, \code{post\_state\_hash} and the radix-trie path. The
  client checks signatures against a bonded-set checkpoint shipped with the browser
  and advanced only through finalized blocks it has verified, requires signatures
  from validators holding more than two thirds of bonded stake, and recomputes the
  path's BLAKE2b-256 hashes up to \code{post\_state\_hash}.
\item[replayed] fetch the process by hash and its inputs at the proof rung, and run
  it locally in \code{Replay}; available only once Obligation~\ref{ob2:replay} is
  discharged.
\end{description}

\subsection{Page verbs}
\code{("lookup", uri, ret)}, \code{("read", name, ret)} (data at a name the site
names by URN or registry entry), \code{("explore", knf\_hash, args, ret)} (exploratory
deploy), \code{("deploy", knf\_hash, args, opts, ret)}, \code{("watch", uri, ch)} and
\code{("session", uri, ret)}. \code{deploy} always prompts, showing the cost from
\code{/api/estimate-cost}, and replies once with the deploy id and again with the
finalization status from \code{/api/deploy-finalization-status}. Pages send code by
hash: the bridge refuses to sign a term the page assembled from strings.

\subsection{Blobs}\label{sec:blobs}
\begin{lstlisting}[language=RustLite]
pub trait BlobSource { fn get(&self, h: [u8; 32]) -> BlobFuture; fn name(&self) -> &str; }
\end{lstlisting}
Implementations, tried in order: the local content cache; mirrors named in the site
manifest and in settings (\code{GET <mirror><hex>}); and, for blobs up to 256\,KiB,
the on-chain layout F1R3Drive already uses --- a metadata map with a first chunk
inline and further chunks on sub-channels --- read through data-at-name. Every source
is untrusted; the hash decides.

\subsection{Sessions}\label{sec:sessions}
\code{("session", uri, ret)} makes the bridge generate a fresh key pair, deploy
\code{svc!("open", pubkey, ack)} signed with it, and reply to the page with a
session name --- an aperture whose sends the bridge turns into deploys (or session-host
messages) signed by that key. The page holds the name; the bridge holds the key.
The service keys its session channel by the public key and accepts a message only
from a deploy whose authenticated deployer (\code{rho:deploy:data},
\code{rho:system:deployerId:ops}, as \code{VersionedRegistry} already does) matches.
\code{DeployData}'s timestamp, \code{validAfterBlockNumber} and
\code{expiration\_timestamp} stop replay. Closing the tab destroys the key.

Interactive traffic goes to a session host: an \code{f1r3x-serve} beside a node,
attached with U7's authenticated handshake, running \code{Deterministic}, which
deploys its log hash to the shard every $N$ rounds or $T$ seconds, whichever first.

%======================================================================
\section{The graded resolver}\label{sec:graded}

\code{gaze-graded} implements \code{Resolver} and the \code{weigh} packing for four
semirings. All arithmetic is on integers.

\begin{center}\small
\begin{tabularx}{\linewidth}{@{}p{2.1cm}p{3.4cm}p{2.8cm}X@{}}
\toprule
Semiring & Representation & $\oplus$, $\otimes$ & Resolver\\
\midrule
Boolean & \code{bool} & or, and & first enabled (today's behaviour)\\
Viterbi & Q32.32 in \code{u64}, $[0,1]$ & max, product rounded down & argmax; ties by candidate order\\
tropical & \code{i64}, $+\infty$ sentinel & min, saturating $+$ & argmin; ties by candidate order\\
$\mathbb{R}_{\geq0}$ & Q32.32 in \code{u64} & saturating $+$, product & sample $\propto$ weight with a 128-bit draw from a xoshiro256** stream seeded from the minter seed\\
\bottomrule
\end{tabularx}
\end{center}

\noindent
In a graded page, a guard is a clause: a Boolean subexpression is crisp ($1$ or
$0$), \code{w(n, d)} is the scalar $n/d$ rounded down to Q32.32, \code{cost(n)} is a
tropical scalar, and \code{and} and \code{or} are $\otimes$ and $\oplus$ --- the
fuzzyware papers' reading of conjunction and disjunction. With \code{"fair": true}
the resolver adds the least positive value $\varepsilon$ to every enabled candidate
before choosing, which is the $\psi\oplus\varepsilon$ weak-fairness combinator.

\begin{requirement}[Ceiling]
The manifest's \code{ceiling} bounds the strength of every clause the tab may run:
\code{crisp} admits Boolean clauses only, \code{idempotent} adds Viterbi and tropical,
\code{sampled} adds $\mathbb{R}_{\geq0}$. The static check covers the body at load;
\code{PageMeter::observe\_consume} checks every continuation at the cut, which covers
code that arrived by reflection and was dropped, and aborts one that exceeds the
ceiling.
\end{requirement}

\noindent
Learning in the fibre, and $V=\mathbb{C}$, are out of scope; the resolver seam is
where they would attach.

%======================================================================
\section{Process model}\label{sec:process}

\begin{description}[leftmargin=1.5em,style=nextline]
\item[P1--P3: one process, a worker per tab.] The shell's main thread owns windows
  and every tab's \code{RhoDocument}. Each tab's \code{TabExec} runs on its own
  worker thread. Events go from \code{RhoEventHandler} to the worker over a channel;
  commit batches come back and are applied in \code{RhoDocument::poll}, which then
  requests a redraw. Synchronous \code{decide} drains block the main thread for at
  most the \code{sync} budget. Isolation between pages is by capability discipline
  and safe Rust.
\item[P4: page processes.] A page process owns its \code{BaseDocument}, layout and
  painting, and ships frames to the browser process as \code{anyrender\_serialize}
  scenes over shared memory; the browser process composites and forwards input. The
  measurement that decides between scene transport and shared GPU textures is
  Obligation~\ref{ob2:ipc}. Process per site, as the proposal requires.
\end{description}

%======================================================================
\section{Devtools, reach, legacy and gateway}\label{sec:devtools}

\paragraph{Devtools.} A \code{rho:gaze:debug} capability, granted only to the
built-in devtools page, attaches to a target tab through \code{f1r3x-host}'s session
vocabulary (\code{pause}, \code{step}, \code{play}, \code{scene}, \code{log}) plus
\code{checkpoints} and \code{revert}. Time travel keeps a ring of soft checkpoints
every 60 frames (default 32 checkpoints); stepping back reverts to the nearest
checkpoint and replays the log to the target frame. The devtools page is written in
\fl{}, which makes it the first real application of the platform. A Rhobots headset
can attach to a tab the same way.

\paragraph{Reach tier.} \code{gaze-reach} compiles the executive, \code{gaze-exec}
and \code{gaze-dom-core} into one \code{cdylib} with the same export conventions as
\code{f1r3x-wasm} (no \code{wasm-bindgen}). \code{gaze-reach.js}, a few hundred
lines, supplies the \code{DomBackend} imports over a table of host elements indexed
by \code{u32}, the frame callback from \code{requestAnimationFrame}, \code{fetch}
with hash verification for \code{net}, and \code{IndexedDB} for \code{store}.
Sessions and deploys are unavailable in the reach tier in v1, since a stock browser
offers no key custody the page cannot reach.

\paragraph{Legacy tab (P4).}\label{sec:legacy} \code{gaze-legacy} embeds the \code{servo} crate at its
current LTS branch in a separate process, connected to the browser process by one
channel carrying ground values encoded as K1G. It gets no grants. CEF behind the
same channel is the fallback.

\paragraph{Gateway (P4).} \code{gaze-gateway} runs \code{RhoDocument} without a
window, runs the page for a bounded number of frames or to quiescence, serialises the
document, and responds with \code{ETag} set to the hash, \code{Content-Digest} using
\code{sha-256} (RFC~9530 registers SHA-2, not BLAKE2b) plus
\code{F1R3-Digest:~blake2b-256=\ldots}, an RFC~9421 signature by the gateway's own
Ed25519 key (secp256k1 is not a registered RFC~9421 algorithm) plus the publisher's
secp256k1 signature over the hash in \code{F1R3-Publisher-Signature}, and
\code{Link:~<f1r3://\ldots>;~rel="alternate"}.

%======================================================================
\section{Testing and conformance}\label{sec:testing}

\begin{description}[leftmargin=1.5em,style=nextline,itemsep=2pt]
\item[Protocol conformance] a corpus of \code{.knf} pages with scripted injections
  and expected committed-DOM hashes per frame, run against \code{gaze-dom-blitz} and
  against \code{gaze-reach} in headless Chromium and Firefox. Both bridges \MUST{}
  pass the same corpus.
\item[Determinism] every corpus run is replayed from its log on \code{x86\_64},
  \code{aarch64} and \code{wasm32}; per-frame \code{batch\_hash} sequences must be
  identical.
\item[Differential] \code{k1ndl1ng-eval} against \code{rho-pure-eval} on generated
  guards over the shared subset; K1G encoding order against \code{models} Score order;
  \Camp{} against \code{rspace++} on recorded logs (Obligation~\ref{ob2:replay}); the
  graded resolvers against \code{fuzzyware/gradedsim.py} on shared examples.
\item[Hostile proxy] a test harness proxy that rewrites, reorders, truncates and
  replays responses between the browser and a local three-validator shard; gate G2
  requires that every rewritten resource is refused and every replayed session
  request rejected.
\item[Fuzzing] \code{cargo-fuzz} targets for the \code{.knf} reader, the K1G decoder,
  the DOM protocol engine, and \code{setHTML} fragments.
\item[Performance] TodoMVC with 1{,}000 items and a 200-mutation interaction: the
  executive and commit portion of a frame \SHOULD{} stay under 4\,ms at p95 on an
  M-series laptop. Missing it triggers the compiled-backend question the bytecode
  note already frames.
\end{description}

%======================================================================
\section{Work packages and schedule}\label{sec:plan}

Phases and gates are the proposal's, from 5 October 2026. The upstream packages
move into P0, because gate G0 --- a program changing rendered text --- cannot pass
without ground strings and apertures. Estimates are engineer-weeks.

\begin{center}\small
\renewcommand{\arraystretch}{1.15}
\begin{tabularx}{\linewidth}{@{}p{1.1cm}p{0.9cm}X p{1.1cm}@{}}
\toprule
WP & Phase & Deliverable & Est.\\
\midrule
U1 & P0 & K1G ground data: AST, parser, normaliser, encoding, matcher, printer & 3\\
U3, U4 & P0 & observer-generic machine; apertures & 2\\
B0 & P0 & \code{RhoDocument} spike: K1G program sets text through a minimal bridge; frame-time and wasm-size measurements & 2\\
E1 & P1 & \code{gaze-knf}, loader, manifest, grounding through a stub broker & 2\\
E2 & P1 & \code{TabExec}, frame loop, recorder, replay & 3\\
D1 & P1 & \code{gaze-dom-core} with the full verb set; \code{gaze-dom-blitz} & 4\\
D2 & P1 & events: \code{RhoEventHandler}, static options, \code{decide} drains & 2\\
T1 & P1 & conformance and determinism suites; TodoMVC in \fl{} & 2\\
C1 & P2 & broker, prompts, revocation; \code{clock}, \code{log}, \code{rand}, \code{net}, \code{store}, \code{nav} & 3\\
S1 & P2 & shard bridge: keystore, resolution, quorum rung, deploys, watches, blobs & 4\\
S2 & P2 & sessions; U7 in \code{f1r3x-serve}; session host with log anchoring & 3\\
N1 & P2 & node: proof endpoints; client light-client check & 4\\
M1 & P2 & page meter; cost quotes & 1\\
U2, U5 & P3 & guards and \code{k1ndl1ng-eval}; resolver seam; \code{Graded} profile & 4\\
G1 & P3 & \code{gaze-graded}: four semirings, fairness, ceiling & 2\\
V1 & P3 & devtools page, time travel, headset attach & 3\\
R1 & P1--P3 & reach tier & 4\\
L1 & P4 & legacy process (Servo; CEF fallback) & 4\\
W1 & P4 & gateway & 2\\
X1 & P4 & page processes and scene transport & 4\\
K1 & P4 & DNS and passkey bootstrap; chrome as a \fl{} page; installers & 4\\
\bottomrule
\end{tabularx}
\end{center}

\noindent
The total is about 62 engineer-weeks over 38 calendar weeks, consistent with the
proposal's three to four engineers. The schedule risk is P0: seven engineer-weeks of
work in four calendar weeks requires the upstream packages to start on day one with
two engineers, one of them familiar with \Camp{}.

%======================================================================
\section{Obligations and decisions}\label{sec:open}

\begin{obligation}\label{ob2:replay}
Build the \code{campf1r3-ispace} adapter and establish whether \Camp{} replays
\code{rspace++} logs; the \emph{replayed} rung depends on it.
\end{obligation}

\begin{obligation}\label{ob2:round}
Specify and implement round-mode resolution over contention components, and state
for which pages the cut-local resolver of Remark~\ref{rem:local} is exact.
\end{obligation}

\begin{obligation}\label{ob2:ipc}
Measure scene serialisation against shared GPU textures for P4, on all three
desktop platforms.
\end{obligation}

\begin{obligation}
Extend the \Kind{} sort-agreement obligation to the K1G tags, and settle whether
K1G's encoding should match the node's \code{Expr} encoding closely enough for
content hashes to agree across the two.
\end{obligation}

\begin{obligation}
Specify the passkey-to-session-key delegation (P-256 WebAuthn assertion authorising
a secp256k1 session key), carried over from the proposal.
\end{obligation}

\subsection{Proposal decisions this specification settles, subject to review}
\begin{description}[leftmargin=1.5em,style=nextline,itemsep=2pt]
\item[DOM granularity] per-element verbs plus structured fragments with
  \code{ref} markers; whole-subtree patches as quoted processes are not needed.
\item[Blob storage] untrusted mirrors first, F1R3Drive's on-chain layout as the
  fallback for small blobs, a node blob endpoint not required.
\item[Repository and name] \code{F1R3FLY-io/f1r3gaze}, working name \Gz{}.
\end{description}

\subsection{Still open}
Compilation of surface \fl{} in the browser (this specification keeps it out of v1);
HTML versus a DDL-defined document language; Servo versus CEF for the legacy tab;
end-to-end encryption versus session-host confinement for private sessions.

%======================================================================
\appendix
\section{The lamp, in K1G}\label{app:lamp}
\begin{lstlisting}
new found, sub, clicks, on, off in {
  doc!("query1", "#lamp", *found) |
  for (@("ok", *lamp) <- found) {
    lamp!("listen", "click", *clicks, {}, *sub) |
    off!(Nil) |
    for (_ <= clicks & _ <- off) { lamp!("classAdd", "lit")    | on!(Nil)  } |
    for (_ <= clicks & _ <- on)  { lamp!("classRemove", "lit") | off!(Nil) }
  }
}
\end{lstlisting}
\noindent
Manifest: \code{imports = [("doc", "rho:gaze:doc")]}, semiring \code{boolean},
ceiling \code{crisp}. The two handlers are persistent continuations that each consume
one click and one token; exactly one token exists, so exactly one fires per click.
The reply to \code{listen} on \code{sub} is left unconsumed, which is permitted.

%======================================================================
\begin{thebibliography}{99}\small
\bibitem{prop} L.~G.~Meredith, \emph{A Browser with f1r3lang as its Native Execution
  Mechanism}, v0.2, \path{publications/f1r3gaze/f1r3-browser-proposal.tex}, 28 Sep 2026.
\bibitem{camp} L.~G.~Meredith, \emph{CampF1R3}, \emph{K1ndl1ng}, \emph{Bytecode, or the
  normal form?}, \path{publications/campf1r3}; implementation in \code{F1R3FLY-io/campf1r3}
  (\code{README.md}, \code{HOSTING.md}).
\bibitem{fuzzy} L.~G.~Meredith, \emph{Graded Where-Clauses} and \emph{Races Decided by the
  Born Rule} (v6), \path{publications/fuzzyware}.
\bibitem{node} \code{F1R3FLY-io/f1r3node-rust}: \path{node/src/rust/web/web_api_routes.rs},
  \path{node/src/rust/web/events_info.rs},
  \path{casper/src/main/resources/VersionedRegistry.rho}, \path{rho-pure-eval},
  \path{rholang/src/rust/interpreter/accounting/costs.rs}.
\bibitem{embers} \code{F1R3FLY-io/embers}: \path{packages/firefly-client},
  \path{packages/embers/src/api}.
\bibitem{drive} \code{F1R3FLY-io/F1R3Drive}: \path{docs/rholang_data_storage.md},
  \path{docs/system_data_flow.md}.
\bibitem{blitz} DioxusLabs, \code{blitz} 0.3.0-beta.2: \path{packages/blitz-dom}
  (\code{Document}, \code{EventHandler}, \code{DocumentMutator}, query selectors),
  \path{packages/blitz-traits/src/net.rs}, \path{packages/blitz-vibey-script}.
\bibitem{rfc9530} R.~Polli, L.~Pardue, \emph{Digest Fields}, RFC~9530, 2024.
\bibitem{rfc9421} A.~Backman, J.~Richer, M.~Sporny, \emph{HTTP Message Signatures},
  RFC~9421, 2024.
\end{thebibliography}
\end{document}
